\documentclass[11pt]{article}
\usepackage[margin=1in]{geometry}
\usepackage{enumitem}
% =============================================================================
% Preambles/header.tex — shared LaTeX/Beamer preamble for this project
%
% Use from a Beamer lecture by prepending:
%     \input{header}
% and compiling with TEXINPUTS=../Preambles:$TEXINPUTS (the /compile-latex
% skill does this automatically).
%
% PALETTE CONTRACT. Color names defined here MUST match the names in
% Quarto/theme-template.scss so Beamer and Quarto renderings use the same
% palette. The scripts/check-palette-sync.sh script enforces this.
%
% To customize: edit the HEX values below AND the matching $variable in
% Quarto/theme-template.scss, then run scripts/check-palette-sync.sh.
% =============================================================================

% -----------------------------------------------------------------------------
% Engine detection + encoding
%
% Our /compile-latex skill uses XeLaTeX, where inputenc is unnecessary and
% can produce encoding warnings. Only load inputenc under pdfTeX.
% -----------------------------------------------------------------------------
\usepackage{iftex}
\ifPDFTeX
  \usepackage[utf8]{inputenc}
\fi

\usepackage{amsmath, amssymb, amsthm}
\usepackage{booktabs}
\usepackage{graphicx}

% -----------------------------------------------------------------------------
% PALETTE — mirrors Quarto/theme-template.scss variable names.
% Load xcolor and define colors BEFORE anything that references them
% (notably hyperref, hypersetup, and the Beamer color assignments).
% -----------------------------------------------------------------------------
\usepackage{xcolor}

\definecolor{primary-blue}{HTML}{012169}      % headings, accents
\definecolor{primary-gold}{HTML}{B9975B}      % emphasis, borders
\definecolor{highlight-yellow}{HTML}{F2A900}  % markers, alerts
\definecolor{light-bg}{HTML}{E8EDF5}          % title / section backgrounds
\definecolor{jet}{HTML}{1A1A1A}               % body text

% Semantic colors (used in diagrams and annotations)
\definecolor{positive}{HTML}{15803D}          % good / observed / identified
\definecolor{negative}{HTML}{B91C1C}          % bad / problematic / bias
\definecolor{neutral}{HTML}{525252}           % reference / context

% Highlight accents (match .hi-* classes in the SCSS)
\definecolor{hi-slate}{HTML}{314F4F}
\definecolor{hi-green}{HTML}{2E8B57}
\definecolor{hi-red}{HTML}{C41E3A}

% Semantic aliases so skill templates and snippets can write
% \textcolor{accent}{...} without hard-coding "primary-blue".
\colorlet{accent}{primary-blue}
\colorlet{accent2}{primary-gold}

% -----------------------------------------------------------------------------
% Hyperref — loaded AFTER xcolor + palette so link colors resolve.
% -----------------------------------------------------------------------------
\usepackage{hyperref}
\hypersetup{colorlinks=true, linkcolor=primary-blue, urlcolor=primary-blue,
            citecolor=primary-blue, pdfborder={0 0 0}}

% -----------------------------------------------------------------------------
% Course code — set once per deck (after \input{header}, before \begin{document})
% via \coursecode{CS401}. Shown in the footer on every slide so a deck's course
% is identifiable without opening the title page. Empty by default.
% -----------------------------------------------------------------------------
\makeatletter
\newcommand{\coursecode}[1]{\gdef\@coursecodetext{#1}}
\gdef\@coursecodetext{}
\makeatother

% -----------------------------------------------------------------------------
% Beamer theme assignments (only apply under Beamer).
%
% We detect Beamer via \@ifclassloaded{beamer}, which is the documented,
% non-internal way. This must be wrapped in \makeatletter/\makeatother
% because \@ifclassloaded contains an @-catcode character.
% -----------------------------------------------------------------------------
\makeatletter
\@ifclassloaded{beamer}{%
  \usefonttheme{professionalfonts}
  \setbeamercolor{structure}{fg=primary-blue}
  \setbeamercolor{frametitle}{fg=primary-blue}
  \setbeamercolor{title}{fg=primary-blue}
  \setbeamercolor{subtitle}{fg=primary-gold}
  \setbeamercolor{author}{fg=primary-blue}
  \setbeamercolor{institute}{fg=neutral}
  \setbeamercolor{date}{fg=primary-gold}
  \setbeamercolor{itemize item}{fg=primary-blue}
  \setbeamercolor{itemize subitem}{fg=primary-gold}
  \setbeamercolor{alerted text}{fg=primary-gold}
  \setbeamercolor{block title}{fg=white, bg=primary-blue}
  \setbeamercolor{block body}{bg=light-bg}
  % Semantic block variants: block=definition/concept (blue), exampleblock=
  % worked example (green/positive), alertblock=key takeaway or caution (gold).
  % Keep to <=2 colored boxes per slide across ALL three variants (INV-7).
  \setbeamercolor{block title example}{fg=white, bg=positive}
  \setbeamercolor{block body example}{bg=light-bg}
  \setbeamercolor{block title alerted}{fg=white, bg=primary-gold}
  \setbeamercolor{block body alerted}{bg=light-bg}

  % Minimal footer: course code (left) + page X / N (right), no navigation clutter
  \setbeamertemplate{navigation symbols}{}
  \setbeamertemplate{footline}{%
    \color{neutral}\footnotesize\hspace{0.7em}\@coursecodetext\hfill
    \insertframenumber/\inserttotalframenumber\hspace{0.7em}\vspace{0.3em}}
}{}
\makeatother

% -----------------------------------------------------------------------------
% TikZ setup — libraries the snippet gallery uses
% -----------------------------------------------------------------------------
\usepackage{tikz}
\usetikzlibrary{arrows.meta, positioning, calc, decorations.pathreplacing,
                fit, shapes.geometric, backgrounds}

% Shared TikZ styles — snippets and hand-written diagrams can pick these up
% via `\begin{tikzpicture}[every node/.style={...}]` or by naming specific
% styles (e.g., `\node[dag-node] (X) at (0,0) {X};`).
\tikzset{
  % Node styles for causal DAGs and similar
  dag-node/.style = {
    draw=primary-blue, fill=light-bg, rounded corners=2pt,
    minimum width=1.2cm, minimum height=0.9cm, align=center, font=\small
  },
  decision-node/.style = {
    draw=primary-gold, fill=white, diamond, aspect=1.8,
    minimum width=1.8cm, minimum height=1.2cm, align=center, font=\small,
    inner sep=1pt
  },
  % Edge styles — observed vs counterfactual
  observed-edge/.style = {-{Latex[length=2.5mm]}, thick, draw=jet},
  counterfactual-edge/.style = {-{Latex[length=2.5mm]}, thick, draw=neutral, dashed},
  confound-edge/.style = {-{Latex[length=2.5mm]}, thick, draw=negative, dashed},
  % Data-point styles for plots
  observed-dot/.style = {fill=primary-blue, draw=primary-blue, circle, inner sep=1.2pt},
  counterfactual-dot/.style = {fill=white, draw=neutral, circle, inner sep=1.2pt}
}

% -----------------------------------------------------------------------------
% Convenience macros
% -----------------------------------------------------------------------------
\newcommand{\muted}[1]{\textcolor{neutral}{#1}}
\newcommand{\key}[1]{\textcolor{primary-gold}{\textbf{#1}}}
\newcommand{\good}[1]{\textcolor{positive}{#1}}
\newcommand{\bad}[1]{\textcolor{negative}{#1}}

% Standout frame macro: full-bleed dark background for section transitions.
% Beamer-only (uses \begin{frame}/\setbeamercolor/\usebeamerfont) -- do not
% call from an article-class file (e.g. Notes/<CODE>/*-notes.tex). \muted,
% \key, \good, \bad, and \coursecode above are plain xcolor/gdef and are
% safe to use from either class; verified when Notes/article-class support
% was added.
% Expand: \transitionslide{Your transition title}
\newcommand{\transitionslide}[1]{%
  \begingroup
  \setbeamercolor{background canvas}{bg=primary-blue}
  \begin{frame}[plain]
    \centering
    \vfill
    {\usebeamerfont{title}\color{white}\Large #1\par}
    \vfill
  \end{frame}
  \endgroup
}

% Section divider frame (Beamer-only), an alternative to \transitionslide with a
% darker two-line style: near-black (jet) background, a small white label line
% above a larger gold title, and a thin gold rule along the bottom edge.
% Expand: \sectiondivider{Section 2}{Pointers}
\newcommand{\sectiondivider}[2]{%
  \begingroup
  \setbeamercolor{background canvas}{bg=jet}
  \begin{frame}[plain]
    \vspace*{3.0cm}
    {\color{white}\ttfamily\large #1\par}
    \vspace{0.5cm}
    {\color{primary-gold}\LARGE #2\par}
    \begin{tikzpicture}[remember picture, overlay]
      \draw[primary-gold, line width=2.5pt]
        ($(current page.south west)+(0,0.1)$) -- ($(current page.south east)+(0,0.1)$);
    \end{tikzpicture}
  \end{frame}
  \endgroup
}

% -----------------------------------------------------------------------------
% Channel watermark — "Concepts That Click" (YouTube).
%
% Article-class files (CompetitiveExam Q/A sets, Notes, Assignments) get a
% light diagonal draft watermark on every page plus a centered footer naming
% the channel. Beamer decks get a subtle TikZ background watermark instead
% (the footline already carries the course code). Centralized here so every
% MCQ PDF is branded without per-file edits.
% -----------------------------------------------------------------------------
\makeatletter
\@ifclassloaded{article}{%
  \usepackage{draftwatermark}
  \SetWatermarkText{Concepts That Click}
  \SetWatermarkScale{1}
  \SetWatermarkFontSize{2cm}
  \SetWatermarkLightness{0.86}
  \SetWatermarkAngle{45}
  \usepackage{fancyhdr}
  \pagestyle{fancy}
  \fancyhf{}
  \fancyfoot[C]{\footnotesize\textcolor{neutral}{Concepts That Click~|~YouTube: \href{https://www.youtube.com/@concepts.thatclick}{@concepts.thatclick}}}
  \renewcommand{\headrulewidth}{0pt}
  \renewcommand{\footrulewidth}{0pt}
  \fancypagestyle{plain}{%
    \fancyhf{}%
    \fancyfoot[C]{\footnotesize\textcolor{neutral}{Concepts That Click~|~YouTube: \href{https://www.youtube.com/@concepts.thatclick}{@concepts.thatclick}}}%
    \renewcommand{\headrulewidth}{0pt}%
    \renewcommand{\footrulewidth}{0pt}%
  }
}{}
\@ifclassloaded{beamer}{%
  \setbeamertemplate{background}{%
    \begin{tikzpicture}[remember picture, overlay]
      \node[rotate=45, opacity=0.055, align=center]
        at (current page.center)
        {\Huge\textcolor{neutral}{Concepts That Click}};
    \end{tikzpicture}%
  }
}{}
\makeatother 
\coursecode{JKSSB}
\renewcommand{\thesection}{24.\arabic{section}}
\newtheorem{example}{Example}[section]
\newenvironment{definitionbox}[1]{%
  \par\noindent\textbf{Definition (#1).}\ \itshape
}{\par\vspace{0.3em}}
\title{Excel Structure --- Detailed Lecture Notes}
\author{JKSSB Computer Awareness}
\date{\today}

\begin{document}
\maketitle

\section{What Excel is and what a workbook is}
Microsoft Excel is a spreadsheet program. A spreadsheet is a grid of boxes
in which you can store text, numbers, dates, and formulas, and in which the
program can calculate results automatically. Excel is one of the standard
office applications, and its saved document is called a \textbf{workbook}.

A \textbf{workbook} is the Excel \emph{file} --- the single item that you
open, edit, and save with a file name. When you start Excel and see a grid,
you are looking at one \textbf{worksheet} that belongs to a workbook. A
workbook can contain many worksheets, and each worksheet is a separate page
of the same file.

\begin{definitionbox}{Workbook}
The Excel file itself. A single workbook can hold one or more worksheets and
is saved under one file name.
\end{definitionbox}

\begin{definitionbox}{Worksheet}
A single sheet of rows and columns inside a workbook. It is also called a
sheet, and it has a tab at the bottom of the Excel window.
\end{definitionbox}

At the bottom of the Excel window you will see one or more \textbf{sheet
tabs} (for example, Sheet1, Sheet2, Sheet3). A new classic workbook opens
with three worksheets, and you can add or remove sheets. Clicking a tab
brings that worksheet to the front. The key point for the exam is the
relationship: the workbook is the container (the file), and the worksheet is
the content (one page inside it).

\section{Workbook vs worksheet}
These two words are a classic near-neighbour pair, so keep the contrast
sharp. The workbook is the file; the worksheet is a part of the file. One
workbook can contain many worksheets, but a worksheet never contains a
workbook. A worksheet is not a separate saved file --- it lives inside the
workbook and is saved as part of it.

\begin{center}
\begin{tabular}{p{0.20\textwidth}p{0.34\textwidth}p{0.34\textwidth}}
\toprule
\textbf{Point} & \textbf{Workbook} & \textbf{Worksheet} \\
\midrule
What it is & The Excel file & One sheet inside a workbook \\
How many & One file & Many per workbook \\
Saved as & A single file name & Part of that file \\
Tab & --- & A tab at the bottom \\
\bottomrule
\end{tabular}
\end{center}

\begin{example}
An item states, ``A worksheet is the Excel file, and a workbook is a part of
it.'' This is wrong. The relation is the other way round: the workbook is
the file, and the worksheet is one page inside it. The wrong option has
swapped the two terms.
\end{example}

\section{The worksheet grid: rows and columns}
Every worksheet is a large grid made of \textbf{rows} and \textbf{columns}.
Rows run horizontally across the sheet and are identified by
\textbf{numbers}, which appear down the left edge (1, 2, 3, \dots).
Columns run vertically down the sheet and are identified by
\textbf{letters}, which appear across the top (A, B, C, \dots).

The column labels do not stop at Z. After A--Z, the labels continue as
two-letter names: AA, AB, \dots, AZ, then BA, BB, \dots, BZ, and so on. This
scheme continues with three-letter names until the final column. In Excel
2007 and later, a worksheet has:

\begin{center}
\begin{tabular}{ll}
\toprule
\textbf{Property} & \textbf{Value in Excel 2007 and later} \\
\midrule
Total rows & 1,048,576 (numbered 1 to 1,048,576) \\
Total columns & 16,384 (labelled A to XFD) \\
First column & A \\
Last column & XFD \\
\bottomrule
\end{tabular}
\end{center}

These two figures are worth memorising because they are common direct-recall
items. In older Excel versions (Excel 97--2003, the \texttt{.xls} format) the
limits were much smaller: 65,536 rows and 256 columns, with the last column
labelled IV. Modern Excel uses the larger grid ending at column XFD.

\section{The cell and the active cell}
A \textbf{cell} is the small rectangular box formed where a row and a column
meet. It is the basic unit of a worksheet, and it is where you actually type
your data. The whole sheet is just a very large collection of cells.

At any moment, exactly one cell is the \textbf{active cell}. The active cell
is the cell that is currently selected: it has a thick border around it, and
anything you type is entered into it. The active cell's address is shown in
the Name Box, and its contents are shown in the Formula Bar. Even when you
select a group of cells, only one of them --- the active cell --- is the one
the next typed entry goes into.

\begin{definitionbox}{Cell}
The rectangular box formed by the intersection of a row and a column, in
which data is entered.
\end{definitionbox}

\begin{definitionbox}{Active cell}
The cell that is currently selected. It has a thick border, receives the
next typed entry, and has its address shown in the Name Box.
\end{definitionbox}

\section{Cell addresses and references}
Every cell has a unique \textbf{cell address} (also called a cell
reference). The address is written with the \textbf{column letter first},
followed by the \textbf{row number}. For example:

\begin{itemize}
  \item \textbf{B5} means column B, row 5.
  \item \textbf{C7} means column C, row 7.
  \item \textbf{D12} means column D, row 12.
\end{itemize}

The order is fixed: letter then number. Writing 12D, B:5, or B\_5 is not a
valid Excel address. Because column labels can have more than one letter, an
address such as AB100 is also valid: it means column AB, row 100.

\section{Ranges}
A \textbf{range} is a group of two or more cells treated as one block. It is
written by typing the addresses of two diagonally opposite corners and
separating them with a \textbf{colon} (:). For example, \textbf{A1:C10}
means all the cells from column A to column C and from row 1 to row 10. A
range is normally rectangular, and the cells in it are contiguous.

The number of cells in a range is the number of columns multiplied by the
number of rows. For example, the range B2:D5 covers three columns (B, C, D)
and four rows (2, 3, 4, 5), so it contains $3 \times 4 = 12$ cells.

\begin{definitionbox}{Range}
A rectangular group of two or more selected cells, written as two corner
addresses joined by a colon, such as A1:C10.
\end{definitionbox}

\begin{example}
A question lists \texttt{B5} and calls it a range. This is wrong: \texttt{B5}
is a \emph{single} cell, because it has only one address and no colon. A
range must have two corner addresses joined by a colon, such as
\texttt{B5:D9}. The colon is the signal that a range, not a single cell, is
meant.
\end{example}

\section{The formula bar}
The \textbf{Formula Bar} is the long white bar that sits below the Ribbon
and above the worksheet grid. It always shows the \emph{contents} of the
active cell. If the active cell holds plain text or a number, the Formula
Bar shows that text or number. If the active cell holds a formula, the
Formula Bar shows the \emph{formula itself} (for example
\texttt{=A1+A2}), while the cell itself shows the calculated result.

You can also edit the active cell in the Formula Bar by clicking in it. This
is useful when a formula is long. Remember the division of labour: the cell
displays the result, and the Formula Bar displays the underlying formula or
data.

\begin{definitionbox}{Formula Bar}
The bar below the Ribbon that shows, and allows editing of, the contents of
the active cell --- including the formula behind a calculated result.
\end{definitionbox}

\section{The name box}
The \textbf{Name Box} is the small box at the \emph{left end of the Formula
Bar}. It normally shows the \textbf{address of the active cell} (for
example, B5). If a range of cells is selected, it shows the name or the
address range of that selection (for example, A1:C10). You can also type an
address into the Name Box and press Enter to jump directly to that cell.

\begin{definitionbox}{Name Box}
The box at the left of the Formula Bar that shows the address of the active
cell or the name of the selected range.
\end{definitionbox}

Keep the two bars apart. The Name Box (on the left) answers ``which cell
am I in?'' The Formula Bar (the long part) answers ``what is in that
cell?''

\section{Basic data types}
Excel accepts several basic kinds of entry in a cell:

\begin{itemize}
  \item \textbf{Text} (also called a label): letters, words, names, and
    headings. Text is not used in arithmetic.
  \item \textbf{Number}: a numeric value that can be used in calculations.
  \item \textbf{Date and time}: entered and displayed as a date or time, but
    stored internally by Excel as a number.
  \item \textbf{Formula}: an entry that begins with the \textbf{=} sign and
    calculates a result from other cells.
\end{itemize}

Excel also recognises \textbf{logical values} (\texttt{TRUE} and
\texttt{FALSE}) and displays \textbf{error values} such as
\texttt{\#VALUE!} when a calculation cannot be completed. By default, Excel
aligns these data types in a helpful way:

\begin{center}
\begin{tabular}{lll}
\toprule
\textbf{Data type} & \textbf{Default alignment} & \textbf{Example} \\
\midrule
Text & Left & Roll No \\
Number & Right & 4500 \\
Date & Right & 15-08-2026 \\
\bottomrule
\end{tabular}
\end{center}

A quick visual check therefore tells you what Excel thinks the entry is: a
value sitting on the \emph{left} is treated as text, and a value sitting on
the \emph{right} is treated as a number (or a date, which is stored as a
number). This default alignment is a common exam point.

\section{Exam retrieval and common confusions}
Six distinctions prevent most errors on this topic.
\begin{enumerate}
  \item A \textbf{workbook} is the file; a \textbf{worksheet} is one page
    inside it. Do not swap them.
  \item A \textbf{cell} is one box; a \textbf{range} is two or more cells
    joined by a colon.
  \item A cell address is \textbf{column letter first, then row number}
    (B5), never the reverse.
  \item A worksheet in Excel 2007+ has \textbf{1,048,576 rows} and
    \textbf{16,384 columns}; the last column is \textbf{XFD}, not IV (IV was
    the old \texttt{.xls} limit).
  \item The \textbf{Name Box} shows the address; the \textbf{Formula Bar}
    shows the contents.
  \item Only \textbf{one cell is active} at a time, even when a range is
    selected.
\end{enumerate}

\begin{example}
An item asks, ``What is the total number of rows in a modern Excel
worksheet?'' Recall the pair: rows are the larger count, 1,048,576, and
columns are 16,384. If the stem instead asks for the last column, the answer
is XFD; if it asks for the number of cells in a range, multiply columns by
rows.
\end{example}

\subsection*{Exercises}
\begin{enumerate}[label=\arabic*.]
  \item Distinguish a workbook from a worksheet.
  \item What is a cell? How is a cell address written, and give one example.
  \item How many rows and how many columns does a worksheet have in Excel
    2007 and later, and what is the last column label?
  \item What is a range? Write the range that covers columns A to C and rows
    1 to 10, and state how many cells it contains.
  \item State what the Name Box and the Formula Bar each display.
  \item List four basic data types that can be entered in an Excel cell and
    give the default alignment of text and of numbers.
\end{enumerate}

\subsection*{Solutions}
\begingroup\small
\begin{enumerate}[label=\arabic*.]
  \item A workbook is the Excel file itself, saved under one file name; a
    worksheet is a single sheet of rows and columns inside a workbook. A
    workbook can contain many worksheets.
  \item A cell is the box formed where a row and a column intersect. A cell
    address is written as the column letter followed by the row number, for
    example B5 (column B, row 5).
  \item A worksheet has 1,048,576 rows and 16,384 columns; the last column is
    labelled XFD.
  \item A range is a rectangular group of two or more cells written as two
    corner addresses joined by a colon. The range is A1:C10, and it contains
    $3 \times 10 = 30$ cells.
  \item The Name Box shows the address of the active cell (or the name of the
    selected range). The Formula Bar shows the contents of the active cell,
    including the formula behind a calculated result.
  \item Four basic data types are text, number, date/time, and formula. Text
    is aligned to the left by default, and numbers to the right by default.
\end{enumerate}
\endgroup
\end{document}
